% GENERATED FILE. Edit canonical lessons, then run npm run generate:books.
% canonical-source-hash: fnv1a64:cd5bd74a29a7c0c6
% canonical-lessons: SA-C345-pratyayah, SA-C345-saksyam, SA-C345-tathyam, SA-C345-sadhayati, SA-C345-niscayah

\chapter{Evidence}
\label{ch:sa-evidence}
\hlchaptermodality{345}

\begin{chapteropening}
\textbf{By the end of this chapter:} \emph{I can say what convinced me and what the evidence shows.}

Complete the last lesson of chapter 345: I can say what convinced me and what the evidence shows.
\end{chapteropening}

\section[\texorpdfstring{\sk{प्रत्ययः} (pratyayaḥ)}{प्रत्ययः (pratyayaḥ)}]{\sk{प्रत्ययः} (pratyayaḥ) — conviction; firm belief}

\label{lesson:SA-C345-pratyayah}



\begin{quote}
Before the new one: say the Sanskrit for consents, then the Sanskrit for is it not so.
\end{quote}

\subsection*{You'll want to know: \sk{प्रत्ययः}}
\textbf{\sk{प्रत्ययः}} — \emph{pratyayaḥ} — \textquotedblleft{}conviction; firm belief\textquotedblright{}.

\begin{quote}
\sk{मम} \sk{प्रत्ययः} \sk{अस्ति} \sk{यत्} \sk{सः} \sk{सत्यं} \sk{वदति।}
\end{quote}

\begin{quote}
\emph{mama pratyayaḥ asti yat saḥ satyaṃ vadati.}
\end{quote}

\begin{quote}
I am convinced that he is telling the truth.
\end{quote}

A masculine noun.

\begin{cousinweb}
\textbf{\sk{प्रत्ययः}} (\emph{pratyayaḥ}) is built on \emph{prati}, \textquotedblleft{}towards\textquotedblright{}, and the root √i, \textquotedblleft{}to go\textquotedblright{}; MW gives \emph{praty-aya} \textquotedblleft{}belief, firm conviction, trust\textquotedblright{} and \textquotedblleft{}proof, ascertainment\textquotedblright{} (MW).
\end{cousinweb}

\begin{grammarlens}[title={What you've built}]
One more word for giving reasons and opinions.
\end{grammarlens}

\subsection*{Your turn}
\begin{itemize}
\raggedright
  \item \emph{Say it:} \emph{pratyayaḥ}
  \item \emph{Say it:} \emph{pratyayaḥ}, once more
  \item \emph{Say it:} \emph{mama pratyayaḥ asti yat saḥ satyaṃ vadati.} — I am convinced that he is telling the truth.
  \item \emph{Recall:} say the Sanskrit for consents, then the Sanskrit for is it not so, then say \emph{pratyayaḥ} again
\end{itemize}

\begin{tcolorbox}[breakable,colback=teal!4,colframe=teal!35!black,title={Before you move on}]
What does \sk{प्रत्ययः} mean? (\textquotedblleft{}Conviction; firm belief\textquotedblright{}.) Say it once more.
\end{tcolorbox}
\section[\texorpdfstring{\sk{साक्ष्यम्} (sākṣyam)}{साक्ष्यम् (sākṣyam)}]{\sk{साक्ष्यम्} (sākṣyam) — testimony; evidence}

\label{lesson:SA-C345-saksyam}



\begin{quote}
Before the new one: say the Sanskrit for is it not so, then the Sanskrit for conviction.
\end{quote}

\subsection*{You'll want to know: \sk{साक्ष्यम्}}
\textbf{\sk{साक्ष्यम्}} — \emph{sākṣyam} — \textquotedblleft{}testimony; evidence\textquotedblright{}.

\begin{quote}
\sk{न्यायाधीशः} \sk{साक्ष्यं} \sk{शृणोति।}
\end{quote}

\begin{quote}
\emph{nyāyādhīśaḥ sākṣyaṃ śṛṇoti.}
\end{quote}

\begin{quote}
The judge hears the testimony.
\end{quote}

A neuter noun.

\begin{cousinweb}
\textbf{\sk{साक्ष्यम्}} (\emph{sākṣyam}) is built on \emph{sākṣin}, \textquotedblleft{}an eyewitness\textquotedblright{}; MW gives \emph{sākṣya} \textquotedblleft{}testimony, evidence, attestation\textquotedblright{} (MW).
\end{cousinweb}

\begin{grammarlens}[title={What you've built}]
One more word for giving reasons and opinions.
\end{grammarlens}

\subsection*{Your turn}
\begin{itemize}
\raggedright
  \item \emph{Say it:} \emph{sākṣyam}
  \item \emph{Say it:} \emph{sākṣyam}, once more
  \item \emph{Say it:} \emph{nyāyādhīśaḥ sākṣyaṃ śṛṇoti.} — The judge hears the testimony.
  \item \emph{Recall:} say the Sanskrit for is it not so, then the Sanskrit for conviction, then say \emph{sākṣyam} again
\end{itemize}

\begin{tcolorbox}[breakable,colback=teal!4,colframe=teal!35!black,title={Before you move on}]
What does \sk{साक्ष्यम्} mean? (\textquotedblleft{}Testimony; evidence\textquotedblright{}.) Say it once more.
\end{tcolorbox}
\section[\texorpdfstring{\sk{तथ्यम्} (tathyam)}{तथ्यम् (tathyam)}]{\sk{तथ्यम्} (tathyam) — a fact; the truth of a matter}

\label{lesson:SA-C345-tathyam}



\begin{quote}
Before the new one: say the Sanskrit for conviction, then the Sanskrit for testimony.
\end{quote}

\subsection*{You'll want to know: \sk{तथ्यम्}}
\textbf{\sk{तथ्यम्}} — \emph{tathyam} — \textquotedblleft{}a fact; the truth of a matter\textquotedblright{}.

\begin{quote}
\sk{एतत्} \sk{तथ्यम्} \sk{अस्ति}, \sk{न} \sk{मिथ्या।}
\end{quote}

\begin{quote}
\emph{etat tathyam asti, na mithyā.}
\end{quote}

\begin{quote}
This is a fact, not a falsehood.
\end{quote}

A neuter noun.

\begin{cousinweb}
\textbf{\sk{तथ्यम्}} (\emph{tathyam}) MW gives \emph{tathya} \textquotedblleft{}being really so, true\textquotedblright{} and the neuter \textquotedblleft{}truth\textquotedblright{}; \textquotedblleft{}a fact\textquotedblright{} is the sense modern Sanskrit gives it, and MW gives no derivation, so its origin is uncertain (MW).
\end{cousinweb}

\begin{grammarlens}[title={What you've built}]
One more word for giving reasons and opinions.
\end{grammarlens}

\subsection*{Your turn}
\begin{itemize}
\raggedright
  \item \emph{Say it:} \emph{tathyam}
  \item \emph{Say it:} \emph{tathyam}, once more
  \item \emph{Say it:} \emph{etat tathyam asti, na mithyā.} — This is a fact, not a falsehood.
  \item \emph{Recall:} say the Sanskrit for conviction, then the Sanskrit for testimony, then say \emph{tathyam} again
\end{itemize}

\begin{tcolorbox}[breakable,colback=teal!4,colframe=teal!35!black,title={Before you move on}]
What does \sk{तथ्यम्} mean? (\textquotedblleft{}A fact; the truth of a matter\textquotedblright{}.) Say it once more.
\end{tcolorbox}
\section[\texorpdfstring{\sk{साधयति} (sādhayati)}{साधयति (sādhayati)}]{\sk{साधयति} (sādhayati) — proves; brings about}

\label{lesson:SA-C345-sadhayati}



\begin{quote}
Before the new one: say the Sanskrit for testimony, then the Sanskrit for a fact.
\end{quote}

\subsection*{You'll want to know: \sk{साधयति}}
\textbf{\sk{साधयति}} — \emph{sādhayati} — \textquotedblleft{}proves; brings about\textquotedblright{}.

\begin{quote}
\sk{अरुणः} \sk{प्रमाणेन} \sk{मतं} \sk{साधयति।}
\end{quote}

\begin{quote}
\emph{aruṇaḥ pramāṇena mataṃ sādhayati.}
\end{quote}

\begin{quote}
Arun proves his view with evidence.
\end{quote}

Past (he …): \sk{सः} \sk{तत्} \sk{साधितवान्}, \emph{saḥ tat sādhitavān}.

\begin{cousinweb}
\textbf{\sk{साधयति}} (\emph{sādhayati}) comes from the root √sidh, \textquotedblleft{}to succeed\textquotedblright{}, in its causative form; MW gives the root the sense \textquotedblleft{}to be proved or established\textquotedblright{} and \emph{sādhita} \textquotedblleft{}brought about, accomplished\textquotedblright{} (MW).
\end{cousinweb}

\begin{grammarlens}[title={What you've built}]
One more word for giving reasons and opinions.
\end{grammarlens}

\subsection*{Your turn}
\begin{itemize}
\raggedright
  \item \emph{Say it:} \emph{sādhayati}
  \item \emph{Say it:} \emph{sādhayati}, once more
  \item \emph{Say it:} \emph{aruṇaḥ pramāṇena mataṃ sādhayati.} — Arun proves his view with evidence.
  \item \emph{Recall:} say the Sanskrit for testimony, then the Sanskrit for a fact, then say \emph{sādhayati} again
\end{itemize}

\begin{tcolorbox}[breakable,colback=teal!4,colframe=teal!35!black,title={Before you move on}]
What does \sk{साधयति} mean? (\textquotedblleft{}Proves; brings about\textquotedblright{}.) Say it once more.
\end{tcolorbox}
\section[\texorpdfstring{\sk{निश्चयः} (niścayaḥ)}{निश्चयः (niścayaḥ)}]{\sk{निश्चयः} (niścayaḥ) — certainty}

\label{lesson:SA-C345-niscayah}



\begin{quote}
Before the new one: say the Sanskrit for a fact, then the Sanskrit for proves.
\end{quote}

\subsection*{You'll want to know: \sk{निश्चयः}}
\textbf{\sk{निश्चयः}} — \emph{niścayaḥ} — \textquotedblleft{}certainty\textquotedblright{}.

\begin{quote}
\sk{मम} \sk{निश्चयः} \sk{अस्ति} \sk{यत्} \sk{सः} \sk{श्वः} \sk{गमिष्यति।}
\end{quote}

\begin{quote}
\emph{mama niścayaḥ asti yat saḥ śvaḥ gamiṣyati.}
\end{quote}

\begin{quote}
I am certain that he will go tomorrow.
\end{quote}

A masculine noun.

\begin{cousinweb}
\textbf{\sk{निश्चयः}} (\emph{niścayaḥ}) is built on \emph{niś}, \textquotedblleft{}out\textquotedblright{}, and the root √ci, \textquotedblleft{}to gather\textquotedblright{}; MW gives \emph{niś-caya} \textquotedblleft{}ascertainment, conviction, certainty\textquotedblright{} (MW).
\end{cousinweb}

\begin{grammarlens}[title={What you've built}]
One more word for giving reasons and opinions.
\end{grammarlens}

\subsection*{Your turn}
\begin{itemize}
\raggedright
  \item \emph{Say it:} \emph{niścayaḥ}
  \item \emph{Say it:} \emph{niścayaḥ}, once more
  \item \emph{Say it:} \emph{mama niścayaḥ asti yat saḥ śvaḥ gamiṣyati.} — I am certain that he will go tomorrow.
  \item \emph{Recall:} say the Sanskrit for a fact, then the Sanskrit for proves, then say \emph{niścayaḥ} again
\end{itemize}

\begin{tcolorbox}[breakable,colback=teal!4,colframe=teal!35!black,title={Before you move on}]
What does \sk{निश्चयः} mean? (\textquotedblleft{}Certainty\textquotedblright{}.) Say it once more.
\end{tcolorbox}
